El calculo de la reposicion responde a la necesidad de determinar cuanto inventario de seguridad es conveniente pedir para cada producto en un nuevo pedido, junto a otra informacion relevante como en que nivel de inventario es mas conveniente hacer el pedido, que cantidad se debe pedir y en que fecha debe hacerse el pedido.

Para esta parte, se utilizaron las formulas de gestion de inventraio clasicas para stock de seguridad, punto de reorden y cantidad económica de pedido (EOQ), los cuales tambien se usaron en los ejemplos de referencia \textit{Ejemplo 1 (SKU Demand Forecasting)} \cite{ejemplo1_demandforecast} y \textit{Ejemplo 2 (Inventory Management with LSTM)} \cite{ejemplo2_lstm_inventory}. En este aspecto, la propuesta de valor de Thary es que estas formulas se hacen en base al pronostico futuro que fue generado en la etapa anterior, y no en base a los datos historicos. 

A partir de las predicciones finales de los meses pronosticados se calculan el promedio mensual $\bar{d}$ y su desviación estándar $\sigma_d$.    

\begin{minted}{python}
def calcular_tabla_reorden(
    df_model, df_mensual, producto_comportamiento, pronostico_futuro_df,
    presupuesto_capital_trabajo=None
):
    #...
    stats_pronostico = pronostico_futuro_df.groupby("producto_id")["prediccion_final"].agg(
        demanda_promedio="mean", desviacion_demanda=lambda x: np.std(x)
    )
\end{minted}

\paragraph{Lead time}

El lead time de cada producto representa el tiempo en el que los proovedores entregarian el producto despues de hacer un pedido, en el sistema se expresa en meses para que sea consistente con la unidad de tiempo de las demandas. Para asegurar esto, si el archivo original contiene el lead time en semanas, el sistema convierte los semanas de lead time a meses, dividiendolas entre 4.345. De no tener el lead time, se asume un valor por defecto de 0.5 meses.

\begin{minted}{python}
lead_time_meses = (
    fila_params["lead_time_semanas"] / 4.345
    if "lead_time_semanas" in df_mensual.columns and not pd.isna(fila_params.get("lead_time_semanas"))
    else LEAD_TIME_DEFAULT_MESES
)
\end{minted}

\paragraph{Stock de seguridad y punto de reorden}

El stock de seguridad es esa cantidad de productos extra que se guarda en el inventario para cubrir la variabilidad de la demanda mientras llega un pedido. 

\[
SS = Z \cdot \sigma_d \cdot \sqrt{L}
\qquad\qquad
\]

\begin{minted}{python}
stock_seguridad = SERVICE_LEVEL_Z * desviacion_demanda * np.sqrt(lead_time_meses)
\end{minted}

El \texttt{SERVICE\_LEVEL\_Z} corresponde a un nivel de servicio de 95\% (Z=1.65), que es el nivel de servicio recomendado por la literatura de gestión de inventarios \cite{chopra2020supply}. Este porcentaje hace referencia a que se espera que el stock de seguridad sea suficiente para cubrir la demanda durante el lead time en 95\% de los casos.

El punto de reorden es el nivel de inventario en el que, cuando se llega a este, se debe de hacer un nuevo pedido.

\[
ROP = \bar{d} \cdot L + SS
\qquad\qquad
\]

\begin{minted}{python}
punto_reorden = demanda_promedio * lead_time_meses + stock_seguridad
\end{minted}

En el que:
\begin{itemize}
    \item $\bar{d}$ es la demanda promedio.
    \item $L$ es el lead time.
    \item $SS$ es el stock de seguridad.
\end{itemize}

\paragraph{Cantidad económica de pedido (EOQ)}

E